% Grupo de 8 posições de uma Swiss Table e a busca por Davi, com etiqueta 89.
% Compartilhado entre o roteiro (figuras.tex) e os slides (hash.tex). Os
% valores das etiquetas são inventados para o exemplo.

\tikzset{
  swctrl/.style = {draw=black, fill=black!5, minimum width=1.25cm,
    minimum height=0.6cm, inner sep=0pt, font=\ttfamily\small, line width=0.6pt},
  swslot/.style = {draw=black, fill=white, minimum width=1.25cm,
    minimum height=0.75cm, inner sep=0pt, font=\ttfamily\small, line width=0.6pt},
  swrot/.style = {font=\sffamily\small, anchor=east},
  swidx/.style = {font=\sffamily\scriptsize, text=black!60},
  swnota/.style = {font=\sffamily\scriptsize, align=center},
  swseta/.style = {-{Stealth[length=2.2mm]}, line width=0.8pt}
}

% Marca de posição vazia no byte de controle: dois hífens, sem ligadura.
\newcommand{\vz}{\mbox{-}\mbox{-}}

\newcommand{\swissGrupo}{%
\begin{tikzpicture}
  \foreach \i/\c/\k in {0/12/Hugo, 1/\vz/, 2/89/Elisa, 3/45/Ana,
                        4/\vz/, 5/89/Davi, 6/3/Gabi, 7/\vz/} {
    \node [swidx] at (\i*1.25, 0.75) {\i};
    \node [swctrl] (c\i) at (\i*1.25, 0) {\c};
    \node [swslot] (s\i) at (\i*1.25, -1.0) {\k};
  }
  % Candidatas: etiqueta igual à procurada.
  \node [swctrl, fill=orange!35] at (c2) {89};
  \node [swctrl, fill=orange!35] at (c5) {89};
  \node [swslot, fill=red!12] at (s2) {Elisa};
  \node [swslot, fill=green!25] at (s5) {Davi};

  \node [swrot] at ($(c0.west)+(-0.15,0)$) {bytes de controle};
  \node [swrot] at ($(s0.west)+(-0.15,0)$) {posições};

  % A etiqueta procurada.
  \node [swctrl, fill=orange!35, label={[font=\sffamily\small]left:busca por \texttt{Davi}, etiqueta}]
    (q) at (2.5*1.25+1.25, 2.0) {89};
  \draw [swseta] (q.south) -- (c2.north east);
  \draw [swseta] (q.south) -- (c5.north west);
  \node [swnota, right=0.3cm of q] {compara com os 8 bytes\\de uma vez};

  \node [swnota, below=0.2cm of s2] {etiqueta igual,\\chave diferente};
  \node [swnota, below=0.2cm of s5] {chave\\encontrada};
\end{tikzpicture}}
